``우리 집 정수기 얼음 나온다!''

얼음을 좋아하는 기룡이는 집에 얼음 정수기를 무려 $N$대나 설치해 두고, 집에 놀러 온 친구들에게 얼음 정수기를 자랑했다.
정수기에는 1번부터 $N$번까지 번호가 붙어 있으며, $i$번 정수기에는 얼음이 $a_i$개 들어 있다. 

모든 정수기는 얼음을 한 번 뽑기 시작하면 들어 있는 모든 얼음이 나온다.
그렇기에 얼음을 버려 낭비하지 않도록 전략적으로 정수기를 선택하는 것이 매우 중요하다.
또한 기룡이는 같은 양의 얼음을 뽑는다 하더라도, 최대한 적은 대수의 정수기를 사용해야 전기세를 아낄 수 있다는 사실을 깨달았다.

기룡이는 친구들에게 시원한 음료를 주기 위해 정확히 $K$개의 얼음이 필요하다. 
정확히 $K$개의 얼음을 얻기 위해 얼음을 뽑아야 하는 정수기 대수의 최솟값을 구하자.